\section{Experiments}\label{sec:experiments}
\subsection{Shared setup and evaluation metrics}
All experiments use the fixed partitions in Table~\ref{tab:splits} and treat
class 1 as positive. Representation fitting and classifier training use the
19,824 training reviews. Validation supports model assessment and deep-model
selection. The four classical configurations use the settings described in
Section~\ref{sec:methods}; no systematic search over $C$, embedding dimensions
or alternative pooling is recorded.

Accuracy, binary precision, recall and F1 are defined by
\begin{align}
 \operatorname{Accuracy}&=\frac{TP+TN}{TP+TN+FP+FN}, &
 \operatorname{Precision}&=\frac{TP}{TP+FP},\\
 \operatorname{Recall}&=\frac{TP}{TP+FN}, &
 F_1&=\frac{2TP}{2TP+FP+FN}.
\end{align}
Rows of each confusion matrix denote gold labels, and columns denote predictions.
We also calculate ROC-AUC in our common evaluation. It measures ranking of
positive versus negative reviews and is valid for SVM margins or their monotonic
sigmoid transforms. Calibration metrics are omitted because the SVM scores are
uncalibrated. The reported F1 values are positive-class scores, not macro averages.

\paragraph{Model-selection qualification.}
We tried an initial TF-IDF configuration and revised it after inspecting both
validation and test metrics. Thus classical test scores
are not a fully untouched holdout estimate. DistilBERT uses validation-only
configuration and checkpoint selection, followed by one final test pass. This
protocol difference and unequal tuning effort must accompany the model comparison.

\subsection{Deep-model training configuration}
Each run minimizes cross-entropy over the two logits using AdamW. Only LoRA
parameters and the two head layers enter the optimizer. All ten epochs run;
the checkpoint with the largest validation F1 is restored afterward. This is
checkpoint selection rather than early stopping.

\begin{table}[htbp]\centering\small
\caption{Training setup of the selected DistilBERT + LoRA run.}\label{tab:dlsetup}
\begin{tabularx}{\linewidth}{@{}lY@{}}\toprule
Setting & Value\\\midrule
Input / maximum length & \code{text_clean}; 512 total tokens, head + tail\\
LoRA & Query/value projections; $r=8$, $\alpha=16$, dropout 0.1\\
Optimizer & AdamW; peak learning rate $3\times10^{-4}$; weight decay 0.01\\
Schedule & Linear warm-up over 6\% of steps, then linear decay to zero\\
Batch size / epochs & 32 / 10; 620 steps per epoch, 6,200 steps total\\
Gradient clipping & Global norm 1.0\\
Precision & fp16 autocast and dynamic loss scaling; parameters held in fp32\\
Seed & 42 for Python, NumPy and PyTorch\\
Device & One NVIDIA Tesla T4 GPU on Kaggle\\
Recorded software & PyTorch 2.11.0+cu130; Transformers 5.20.2; PEFT 0.21.2\\
Peak allocated GPU memory & 2,910.54 MiB in the final run\\\bottomrule
\end{tabularx}
\end{table}

\subsection{Validation search and checkpoint selection}\label{sec:tuning}
Eight configurations are evaluated in a one-factor-at-a-time search centered on
$r=8$, query/value adaptation, learning rate $3\times10^{-4}$, length 256 and
head-only truncation. Variations test learning rates $10^{-4}$ and $5\times10^{-4}$;
ranks 4 and 16 with $\alpha=2r$; adaptation of all six linear modules in each
block; length 512 with head-only input; and length 512 with head-plus-tail input.
All other training settings remain fixed. The ``all'' option includes
\code{q_lin}, \code{k_lin}, \code{v_lin}, \code{out_lin}, \code{lin1} and \code{lin2}.

\begin{table}[htbp]\centering\small
\caption{Eight search configurations, ranked by best validation F1. ``qv'' means query/value adaptation. Each row uses its own best epoch.}\label{tab:search}
\begin{tabular}{@{}lrlrrrr@{}}\toprule
Length, strategy & $r$ & Modules & LR & Epoch & Val. acc. & Val. F1\\\midrule
\tableinput{tables/dl_search.tex}
\bottomrule\end{tabular}
\end{table}
\reportfigure{person3_image2.png}{0.94}{Best validation F1 of the eight search runs. The horizontal axis begins at 0.905, so bar lengths emphasize small differences.}{fig:search}

The selected configuration obtains validation F1 0.9301 at epoch 6. Extending
head-only input from 256 to 512 raises F1 from 0.9138 to 0.9279; head-plus-tail
adds approximately 0.0022. Adapting all modules at length 256 raises F1 to 0.9175,
but increases trainable parameters to 1,255,682 and shows overfitting: its
training loss declines to approximately 0.031 while validation loss grows to
0.478 by epoch 10. Increasing adapter rank alone does not improve this search's
best F1. These observations concern one seed and do not test all interactions.

The selected configuration is retrained from the pretrained checkpoint, using
training records only. Its best validation epoch is 8, with F1 0.9298, rather
than search epoch 6. The saved logs distinguish these runs. Numerical
nondeterminism is a plausible contributor to the difference; it is not evidence
that the test set selected the epoch. The final checkpoint predicts validation
and test, saving IDs, labels and positive-class scores. The final script refuses
to overwrite existing results unless explicitly forced.

\subsection{Reproducibility and evidence checks}
The preprocessing notebook prepares the fixed data. Training is implemented in
\code{ml/model.py} for classical models and \code{dl/train_lora.py} for the deep
model. Data preparation records spaCy 3.8.16,
\code{en_core_web_sm} 3.8.0 and dataset export checksums. Word2Vec uses seed 42
and one worker, although numerical environments can still affect reproduction.
We observed a Gensim runtime warning during development. Training completed
with finite vectors, but numerical environments can still affect reproduction.

We evaluated all five models from prediction files and cross-checked the metrics
against our training results to the displayed precision. The saved confusion
counts reproduce accuracy, precision, recall and F1 in \code{results/}.
We also verify the DistilBERT confusion matrix and ROC-AUC from its 25,000 saved
test predictions, category counts from the 100 final annotation records, and
the prepared data checksums.

Our submission includes the standalone evaluation and error-analysis scripts
in \pathcode{.eval_demo_package/eval/}, with their recorded outputs in \pathcode{.eval_demo_package/results/eval/}.
The four classical per-review prediction files are not included; their ROC-AUC
values and shared-error counts are documented in these recorded outputs.
To regenerate the prediction files and repeat the evaluation, use
\pathcode{.eval_demo_package/eval/make_ml_preds.py} and the commands in Appendix~\ref{app:reproduce}.

\subsection{Interactive Streamlit demonstration}\label{sec:demo}
We include the demonstration source in \pathcode{.eval_demo_package/demo/}.
Our app implements a text area, preset examples, and a hard test case with
an optional true-label reveal. It displays all five predicted labels,
positive-class outputs, a comparison bar chart, model-loading status and test
accuracy/F1 in the sidebar. The app handles empty input and indicates unavailable
models. Figures~\ref{fig:demooverview}--\ref{fig:demolong} are the running-app
screenshots from our demonstration.

ML bundles are rebuilt from training using the original settings; the app loads
the committed LoRA adapter and classification head with the pretrained base.
Its inference mirrors classical preprocessing and the selected head-tail input.
The first transformer load requires downloading the base model. Our recorded
fresh-environment test passed 12/12 automated app checks, covering input,
model-loading, scoring and length-warning behavior. Setup and predictor-check
commands are listed in Appendix~\ref{app:reproduce}.

\reportfigure{person4_rId34.png}{0.94}{Streamlit at startup. The sidebar lists the five models as loaded and shows their saved test-set scores.}{fig:demooverview}

\begin{table}[htbp]\centering\small
\caption{Recorded demo outputs. Columns are LR and SVM for TF-IDF and Word2Vec, followed by DistilBERT + LoRA. SVM entries are transformed scores.}\label{tab:demo}
\begin{tabularx}{\linewidth}{@{}Yrrrrr@{}}\toprule
Input (shortened) & TF LR & TF SVM & W2V LR & W2V SVM & DL\\\midrule
``An absolute delight from start to finish''\ldots & .8611 & .7769 & .9980 & .9145 & .9993\\
``It's not too bad.'' & .0013 & .0831 & .0000 & .0046 & .9347\\
``hilarious antics ensue (I wish)'' & .8037 & .7046 & .9996 & .9598 & .9971\\
``stupid and lame for the first half''\ldots ``good for the last half'' & .1227 & .2156 & .0000 & .0048 & .7638\\\bottomrule
\end{tabularx}
\end{table}
The examples expose differences in composition: DistilBERT predicts positive
for the negation and mixed-opinion inputs, while the classical models predict
negative. Every model predicts positive for the sarcastic fragment, illustrating
a shared failure. Short demo fragments are new inputs and need not yield the
same prediction as the full test reviews from which some were drawn.

\reportfigure{person4_rId35.png}{0.94}{Demo prediction for ``It's not too bad.'' The classical models return negative; DistilBERT returns positive.}{fig:demonegation}
\reportfigure{person4_rId36.png}{0.94}{Long-review demonstration for \code{test_14218}. The warning's 512-token limit includes two special tokens; the reported 1,083 input wordpieces also includes them.}{fig:demolong}

The interface names its numeric column \code{P(positive)}, but also warns that
SVM outputs are sigmoid-transformed margins. Confidence magnitudes should not
be compared across the five models as calibrated probabilities. The length
warning indicates that the deep model sees only the beginning and end; it does
not establish that truncation caused a particular error. Run instructions and
artifact locations are provided in Appendix~\ref{app:reproduce}.
\FloatBarrier
